\SourceLeaf{142}

\[
(Ab+B(a+d)+Cc)(ad-bc)=(-Ab+B(a+d)-Cc)e'e'
\]

adeoque

\[
B(a+d)=0
\]

Denique ex $a[3]-b[2]$ fit

\[
(Bb+Cd)(ad-bc)=(-Bb+Ca)e'e'
\]

adeoque

\[
C(a+d)=0
\]

Quare quum omnes $A,B,C$ nequeant esse $=0$, necessario erit $a+d=0$ sive $a=-d$.

Ex $a[2]-b[1]$ fit

\[
(Ba+Cc)(ad-bc)=(Ba-Ab)e'e'
\]

unde

\[
Ab-2Ba-Cc=0\quad\mathinner{.\ .\ .\ .\ .\ .\ .\ .}\quad[5]
\]

Ex aequationibus $e+e'=0$, $a+d=0$ sive

\[
\alpha\delta-\beta\gamma+\alpha'\delta'-\beta'\gamma'=0,\quad\alpha\delta'-\beta\gamma'-\gamma\beta'+\delta\alpha'=0
\]

sequitur\quad $(\alpha+\alpha')(\delta+\delta')=(\beta+\beta')(\gamma+\gamma')$\quad sive

\[
(\alpha+\alpha'):(\gamma+\gamma')=(\beta+\beta'):(\delta+\delta')
\]

Sit rationi huic\HistoricalNote{*)}{Si omnes $\alpha+\alpha'$, $\gamma+\gamma'$, $\beta+\beta'$, $\delta+\delta'$ essent $=0$, ratio indeterminata foret, adeoque methodus non applicabilis. Sed exigua attentio docet, hoc cum suppositionibus nostris consistere non posse. Foret enim $\alpha\delta-\beta\gamma=\alpha'\delta'-\beta'\gamma'$ \emph{i. e.} $e=e'$ adeoque, quia $e=-e'$, $e=e'=0$. Hinc vero etiam $B'B'-A'C'$ \emph{i. e.} determinans formae $F'$ fieret $=0$, quales formas omnino exclusimus.} in numeris minimis aequalis ratio $m:n$, ita ut $m,n$ inter se primi sint, accipianturque $\mu,\nu$ ita ut fiat $\mu m+\nu n=1$. Porro sit $r$ div. comm. max. numerorum $a,b,c$; cuius quadratum propterea metietur ipsum $aa+bc$ sive $bc-ad$ sive $ee$; quare $r$ etiam ipsum $e$ metietur. His ita factis, si forma $F$ per substitutionem

\[
x=mt+\frac{\nu e}{r}u,\qquad y=nt-\frac{\mu e}{r}u
\]

in formam\quad $Mtt+2Ntu+Puu$\quad $(G)$ transire supponitur, haec anceps erit formamque $F'$ implicabit.

\SourceLeaf{143}

\emph{Dem.} I. Quo pateat, formam $G$ esse ancipitem, ostendemus esse

\[
M(b\mu\mu-2a\mu\nu-c\nu\nu)=2Nr
\]

unde quia $r$ ipsos $a,b,c$ metitur $\frac{1}{r}(b\mu\mu-2a\mu\nu-c\nu\nu)$ integer erit, adeoque $2N$ multiplum ipsius $M$. Erit autem

\[
\begin{gathered}M=Amm+2Bmn+Cnn,\\Nr=(Am\nu-B(m\mu-n\nu)-Cn\mu)e\quad\mathinner{.\ .}\quad[6]\end{gathered}
\]

Porro per evolutionem facile confirmatur esse

\[
\begin{aligned}2e+2a&=e-e'+a-d\\&=(\alpha-\alpha')(\delta+\delta')-(\beta-\beta')(\gamma+\gamma')\end{aligned}
\]

\[
2b=(\alpha+\alpha')(\beta-\beta')-(\alpha-\alpha')(\beta+\beta')
\]

Hinc quoniam $m(\gamma+\gamma')=n(\alpha+\alpha')$, $m(\delta+\delta')=n(\beta+\beta')$ erit

\[
m(2e+2a)=-2nb\quad\text{sive}
\]

\[
me+ma+nb=0\quad\mathinner{.\ .\ .\ .\ .\ .\ .\ .\ .}\quad[7]
\]

Eodem modo erit

\[
\begin{aligned}2e-2a&=e-e'-a+d\\&=(\alpha+\alpha')(\delta-\delta')-(\beta+\beta')(\gamma-\gamma')\end{aligned}
\]

\[
2c=(\gamma-\gamma')(\delta+\delta')-(\gamma+\gamma')(\delta-\delta')
\]

atque hinc $n(2e-2a)=-2mc$ sive

\[
ne-na+mc=0\quad\mathinner{.\ .\ .\ .\ .\ .\ .\ .\ .}\quad[8]
\]

Iam si ad $mm(b\mu\mu-2a\mu\nu-c\nu\nu)$ additur

\[
\begin{gathered}(1-m\mu-n\nu)(m\nu(e-a)+(m\mu+1)b)\\+(me+ma+nb)(m\mu\nu+\nu)+(ne-na+mc)m\nu\nu\end{gathered}
\]

quod manifesto $=0$, propter

\[
1-\mu m-\nu u=0,\quad me+ma+nb=0,\quad ne-na+mc=0
\]

prodit productis rite evolutis partibusque se destruentibus deletis, $2m\nu e+b$. Quare erit

\[
mm(b\mu\mu-2a\mu\nu-c\nu\nu)=2m\nu e+b\quad\mathinner{.\ .\ .\ .\ .}\quad[9]
\]

Eodem modo addendo ad $mn(b\mu\mu-2a\mu\nu-c\nu\nu)$ haec:

\SourceLeaf{144}

\[
\begin{gathered}(1-m\mu-n\nu)((n\nu-m\mu)e-(1+m\mu+n\nu)a)\\-(me+ma+nb)m\mu\mu+(ne-na+mc)n\nu\nu\end{gathered}
\]

invenitur

\[
mn(b\mu\mu-2a\mu\nu-c\nu\nu)=(n\nu-m\mu)e-a\quad .\ .\ .\ .\ [10]
\]

Denique addendo ad $nn(b\mu\mu-2a\mu\nu-c\nu\nu)$ haec:

\[
\begin{gathered}(m\mu+n\nu-1)(n\mu(e+a)+(n\nu+1)c)\\-(me+ma+nb)n\mu\mu-(ne-na+mc)(n\mu\nu+\mu)\end{gathered}
\]

fit

\[
nn(b\mu\mu-2a\mu\nu-c\nu\nu)=-2n\mu e-c\quad .\ .\ .\ .\ [11]
\]

Iam ex $9$, $10$. $11$, deducitur

\[
\begin{gathered}(Amm+2Bmn+Cnn)(b\mu\mu-2a\mu\nu-c\nu\nu)\\=2e(Am\nu+B(n\nu-m\mu)-Cn\mu)+Ab-2Ba-Cc\end{gathered}
\]

sive propter $[6],$

\[
M(b\mu\mu-2a\mu\nu-c\nu\nu)=2Nr.\quad Q.\ E.\ D.
\]

II.\quad Ut probetur, formam $G$ implicare formam $F'$, demonstrabimus, \emph{primo} $G$ transire in $F'$ ponendo

\[
\begin{gathered}t=(\mu\alpha+\nu\gamma)x'+(\mu\beta+\nu\delta)y',\\u=\frac re(n\alpha-m\gamma)x'+\frac re(n\beta-m\delta)y'\ .\ .\ .\ (S)\end{gathered}
\]

\emph{secundo} $\frac{r}{e}(n\alpha-m\gamma)$, $\frac{r}{e}(n\beta-m\delta)$ esse integros.

1.\quad Quoniam $F$ transit in $G$ ponendo

\[
x=mt+\frac{\nu e}{r}u,\quad y=nt-\frac{\mu e}{r}u
\]

forma $G$ per substitutionem $(S)$ transmutabitur in eandem formam in quam $F$ transformatur ponendo

\[
\begin{aligned}x={}&m((\mu\alpha+\nu\gamma)x'+(\mu\beta+\nu\delta)y')\\&+\nu((n\alpha-m\gamma)x'+(n\beta-m\delta)y')\end{aligned}
\]

\emph{i. e.}\qquad $=\alpha(m\mu+n\nu)x'+\beta(m\mu+n\nu)y'$\quad sive\quad $=\alpha x'+\beta y'$

et

\[
\begin{aligned}y={}&n((\mu\alpha+\nu\gamma)x'+(\mu\beta+\nu\delta)y')\\&-\mu((n\alpha-m\gamma)x'+(n\beta-m\delta)y')\end{aligned}
\]

\emph{i. e.}\qquad $=\gamma(n\nu+m\mu)x'+\delta(n\nu+m\mu)y'$\quad sive\quad $=\gamma x'+\delta y'$

\SourceLeaf{145}

Per hanc vero substitutionem $F$ transit in $F'$: quare per substitutionem $(S)$ etiam $G$ transibit in $F'$.

2.\quad Ex valoribus ipsorum $e,b,d$ invenitur $\alpha'e+\gamma b-\alpha d=0$, sive propter $d=-a$, $n\alpha'e+n\alpha a+n\gamma b=0$; hinc ex $[7]$, $n\alpha'e+n\alpha a=m\gamma e+m\gamma a$ sive

\[
(n\alpha-m\gamma)a=(m\gamma-n\alpha')e\quad .\ .\ .\ .\ .\ .\ .\ [12]
\]

Porro fit $\alpha nb=-\alpha m(e+a)$, $\gamma mb=-m(\alpha'e+\alpha a)$ adeoque

\[
(n\alpha-m\gamma)b=(\alpha'-\alpha)me\quad .\ .\ .\ .\ .\ .\ .\ [13]
\]

Denique fit $\gamma'e-\gamma a+\alpha c=0$: hinc multiplicando per $n$, et pro $na$ substituendo valorem ex $[8]$ fit

\[
(n\alpha-m\gamma)c=(\gamma-\gamma')ne\quad .\ .\ .\ .\ .\ .\ .\ [14]
\]

Simili modo eruitur $\beta'e+\delta b-\beta d=0$, sive $n\beta'e+n\delta b+n\beta a=0$, adeoque per $[7],$ $n\beta'e+n\beta a=m\delta e+m\delta a$ sive

\[
(n\beta-m\delta)a=(m\delta-n\beta')e\quad .\ .\ .\ .\ .\ .\ .\ [15]
\]

Porro fit $\beta nb=-\beta m(e+a)$, $\delta mb=-m(\beta'e+\beta a)$ adeoque

\[
(n\beta-m\delta)b=(\beta'-\beta)me\quad .\ .\ .\ .\ .\ .\ .\ [16]
\]

Tandem $\delta e-\delta a+\beta c=0$: hinc multiplicando per $n$ et substituendo pro $na$ valorem ex $[8]$ fit

\[
(n\beta-m\delta)c=(\delta-\delta')ne\quad .\ .\ .\ .\ .\ .\ .\ [17]
\]

Iam quum divisor communis maximus numerorum $a,b,c$ sit $r$, integri $\mathfrak{A},\mathfrak{B},\mathfrak{C}$ ita accipi possunt, ut fiat

\[
\mathfrak{A}a+\mathfrak{B}b+\mathfrak{C}c=r
\]

Quo facto erit ex $12$, $13$, $14$; $15$, $16$, $17$

\[
\begin{aligned}&\mathfrak A(m\gamma-n\alpha')+\mathfrak B(\alpha'-\alpha)m+\mathfrak C(\gamma-\gamma')n\\&\qquad=\frac re(n\alpha-m\gamma)\\&\mathfrak A(m\delta-n\beta')+\mathfrak B(\beta'-\beta)m+\mathfrak C(\delta-\delta')n\\&\qquad=\frac re(n\beta-m\delta)\end{aligned}
\]

adeoque $\frac{r}{e}(n\alpha-n\gamma),\quad\frac{r}{e}(n\beta-m\delta)$ integri.\quad $Q.\ E.\ D.$

\SourceLeaf{146}

\Article{165}

\emph{Ex.}\quad Forma $3xx+14xy-4yy$ in formam $-12x'x'-18x'y'+39y'y'$ transmutatur, tum proprie, ponendo

\[
x=4x'+11y',\quad y=-x'-2y'
\]

tum improprie, ponendo

\[
x=-74x'+89y',\quad y=15x'-18y'
\]

Hic igitur $\alpha+\alpha'$, $\beta+\beta'$, $\gamma+\gamma'$, $\delta+\delta'$ sunt $-70$, $100$, $14$, $-20$; est autem $-70:14=100:-20=5:-1$.\quad Faciemus itaque $m=5$, $n=-1$, $\mu=0$, $\nu=-1$.\quad Numeri autem $a,b,c$ inveniuntur $-237$, $-1170$, $48$, quorum divisor communis maximus $=3=r$; denique fit $e=3$.\quad Hinc transformatio $(S)$ haec erit, $x=5t-u$, $y=-t$.\quad Per quam forma $(3,7,-4)$ transit in formam ancipitem $tt-16tu+3uu$.

Si formae $F,F'$ sunt aequivalentes: forma $G$, sub $F$ contenta, etiam sub $F'$ contenta erit.\quad Sed quoniam eandem formam etiam implicat, ipsi aequivalens erit, et proin etiam formae $F$.\quad In hoc igitur casu theorema ita enunciabitur:

\emph{Si $F,F'$ tam proprie, quam improprie sunt aequivalentes: forma anceps utrique aequivalens inveniri poterit.} ---\quad Ceterum in hoc casu $e=\pm1$, adeoque etiam $r.$ ipsum $e$ metiens, $=1$ erit.

Haec de formarum transformatione in genere sufficiant: transimus itaque ad considerationem \emph{repraesentationum}.

\begin{center}\emph{Generalia de repraesentationibus numerorum per formas, earumque nexu cum transformationibus.}\end{center}

\Article{166}

\emph{Si forma $F$ formam $F'$ implicat: quicunque numerus per $F'$ repraesentari potest, etiam per $F$ poterit.}

Sint indeterminatae formarum $F,F'$ respective $x,y;x',y'$, ponamusque numerum $M$ per $F'$ repraesentari faciendo $x'=m$, $y'=n$, formam $F$ vero in $F'$ transire per substitutionem

\[
x=\alpha x'+\beta y',\quad y=\gamma x'+\delta y'
\]

\SourceLeaf{147}

Tum manifestum est, si ponatur

\[
x=\alpha m+\beta n,\quad y=\gamma m+\delta n
\]

$F$ transire in $M$.

Si $M$ pluribus modis per formam $F'$ repraesentari potest, e. g. etiam faciendo $x'=m'$, $y'=n'$: plures repraesentationes ipsius $M$ per $F$ inde sequentur. Si enim esset tum

\[
\alpha m+\beta n=\alpha m'+\beta n'\quad\text{tum}\quad\gamma m+\delta n=\gamma m'+\delta n'
\]

foret aut $\alpha\delta-\beta\gamma=0$, adeoque etiam determinans formae $F=0$ contra hyp., aut $m=m'$, $n=n'$.\quad Hinc sequitur $M$ ad minimum totidem modis diversis per $F$ repraesentari posse quot per $F'$.

Si igitur tum $F$ ipsam $F'$, tum $F'$ ipsam $F$ implicat \emph{i. e.} si $F,F'$ sunt aequivalentes, numerusque $M$ per alterutram repraesentari potest: etiam per alteram repraesentari poterit, et quidem totidem modis diversis per alteram, quot per alteram.

Denique observamus, in hocce casu divisorem communem maximum numerorum $m,n$ aequalem esse divisori comm. max. numerorum $\alpha m+\beta n$, $\gamma m+\delta n$. Sit ille $=\Delta$, numerique $\mu,\nu$ ita accepti, ut fiat $\mu m+\nu n=\Delta$.\quad Tum erit

\[
\begin{gathered}(\delta\mu-\gamma\nu)(\alpha m+\beta n)-(\beta\mu-\alpha\nu)(\gamma m+\delta n)\\=(\alpha\delta-\beta\gamma)(\mu m+\nu n)=\pm\Delta\end{gathered}
\]

Hinc div. comm. max. numerorum $\alpha m+\beta n$, $\gamma m+\delta n$ metietur ipsum $\Delta$, $\Delta$ vero etiam illum metietur, quia manifesto ipsos $\alpha m+\beta n$, $\gamma m+\delta n$ metitur.\quad Quare necessario ille erit $=\Delta.$ ---\quad Quando igitur $m,n$ inter se primi sunt, etiam $\alpha m+\beta n.\quad\gamma m+\delta n$ inter se primi erunt.

\Needspace{12\baselineskip}
\Article{167}

\textsc{Theorema.}\quad \emph{Si formae}

\[
axx+2bxy+cyy\quad .\ .\ .\ .\ .\ .\ (F)
\]

\[
a'x'x'+2b'x'y'+c'y'y'\quad .\ .\ .\ .\ .\ .\ (F')
\]

\emph{sunt aequivalentes, ipsarum determinans $=D$, posteriorque in priorem transit ponendo}

\[
x'=\alpha x+\beta y,\quad y'=\gamma x+\delta y
\]

\SourceLeaf{148}

\emph{porro numerus $M$ per $F$ repraesentatur, faciendo $x=m$, $y=n$, adeoque per $F'$ faciendo}

\[
x'=\alpha m+\beta n=m',\quad y'=\gamma m+\delta n=n'
\]

\emph{et quidem ita ut $m$ ad $n$ eoque ipso etiam $m'$ ad $n'$ sit primus: ambae repraesentationes aut ad eundem valorem expressionis $\surd D(\operatorname{mod.}M)$ pertinebunt, aut ad oppositos, prout transformatio formae $F'$ in $F$ propria est vel impropria.}

\emph{Dem.}\quad Determinentur numeri $\mu,\nu$ ita ut fiat $\mu m+\nu n=1$, ponaturque

\[
\frac{\delta\mu-\gamma\nu}{\alpha\delta-\beta\gamma}=\mu',\quad \frac{-\beta\mu+\alpha\nu}{\alpha\delta-\beta\gamma}=\nu'
\]

(qui erunt integri propter\quad $\alpha\delta-\gamma\beta=\pm1$).\quad Tum erit

\[
\mu'm'+\nu'n'=1.\quad\text{(Cf. art. praec. fin.)}
\]

Porro sit

\[
\begin{gathered}\mu(bm+cn)-\nu(am+bn)=V,\\\mu'(b'm'+c'n')-\nu'(a'm'+b'n')=V'\end{gathered}
\]

eruntque $V,V'$ valores expr. $\surd D(\operatorname{mod.}M)$ ad quos repraesentatio prima et secunda pertinent.\quad Si in $V'$ pro $\mu',\nu',m',n'$ valores ipsorum substituuntur, in $V$ vero

\[
\begin{array}{lll}\text{pro}&a,&a'\alpha\alpha+2b'\alpha\gamma+c'\gamma\gamma\\\text{pro}&b,&a'a\beta+b'(\alpha\delta+\beta\gamma)+c'\gamma\delta\\\text{pro}&c,&a'\beta\beta+2b'\beta\delta+c'\delta\delta\end{array}
\]

invenietur evolutione facta\quad $V=V'(\alpha\delta-\beta\gamma).$

Quare erit aut $V=V'$, aut $V=-V'$, prout $\alpha\delta-\beta\gamma=+1$ aut $=-1$, \emph{i. e.} repraesentationes pertinebunt ad eundem valorem expr. $\surd D(\operatorname{mod.}M)$ vel ad oppositos, prout transformatio formae $F'$ in $F$ est propria vel impropria.\quad $Q.\ E.\ D.$

Si itaque plures repraesentationes numeri $M$ per formam $(a,b,c)$, ope valorum inter se primorum indeterminatarum $x,y$, habentur ad valores \emph{diversos} expr. $\surd D(\operatorname{mod.}M)$ pertinentes: repraesentationes respodentes per formam $(a',b',c')$ ad eosdem resp. valores pertinebunt, et si nulla repraesentatio numeri $M$ per formam aliquam ad valorem quendam determinatum pertinens datur, nulla quoque dabitur ad hunc valorem pertinens per formam illi aequivalentem.

\SourceLeaf{149}

\Article{168}

\textsc{Theorema.}\quad \emph{Si numerus $M$ per formam $axx+2bxy+cyy$ repraesentatur tribuendo ipsis $x,y$, valores inter se primos $m,n$, valorque expressionis $\surd D(\operatorname{mod.}M)$, ad quem haec repraesentatio pertinet, est $N$: formae $(a,b,c)$, $(M,N,\frac{NN-D}{M})$ proprie aequivalentes erunt.}

\emph{Demonstr.}\quad Ex art. $155$ patet, numeros integros $\mu,\nu$ inveniri posse ita ut sit

\[
m\mu+n\nu=1,\quad \mu(bm+cn)-\nu(am+bn)=N
\]

Quo facto forma $(a,b,c)$ per substitutionem $x=mx'-\nu y'$, $y=nx'+\mu y'$, quae manifesto est propria, transit in formam cuius determinans $=$ $D(m\mu+n\nu)^2$ \emph{i. e.} $=D$, sive in formam aequivalentem: quae forma si ponitur $=(M',N',\frac{N'N'-D}{M'})$, erit

\[
\begin{gathered}M'=amm+2bmn+cnn=M,\\N'=-m\nu a+(m\mu-n\nu)b+n\mu c=N\end{gathered}
\]

Quare forma in quam $(a,b,c)$ per transformationem illam mutatur, erit $(M,N,\frac{NN-D}{M})$.\quad $Q.\ E.\ D.$

Ceterum ex aequationibus

\[
m\mu+n\nu=1,\quad \mu(mb+nc)-\nu(ma+nb)=N
\]

deducitur

\[
\begin{gathered}\mu=\frac{nN+ma+nb}{amm+2bmn+cnn}=\frac{nN+ma+nb}{M},\\\nu=\frac{mb+nc-mN}{M}\end{gathered}
\]

qui numeri itaque erunt integri.

Porro observandum, hanc propositionem locum non habere, si $M=0$; tum enim terminus $\frac{NN-D}{M}$ fit indeterminatus\HistoricalNote{*)}{In hoc enim casu, si ad ipsum phrasin extendere volumus, haec: $N$ esse valorem expr. $\surd D(\operatorname{mod.}M)$, sive $NN\equiv D(\operatorname{mod.}M)$ significabit, $NN-D$ esse multiplum ipsius $M$, adeoque $=0$.}.

\Article{169}

Si plures repraesentationes numeri $M$, per $(a,b,c)$ habentur. ad eundem valorem expr. $\surd D(\operatorname{mod.}M)$, $N$, pertinentes (ubi valores ipsorum $x,y$ semper inter se primos supponimus): plures etiam transformationes propriae formae $(a,b,c)..(F)$, in $(M.\ N,\frac{NN-D}{M})..(G)$ inde deducentur.\quad Scilicet si etiam per hos valores $x=m'$, $y=n'$ talis repraesentatio provenit, $(F)$ etiam per substitutionem

\SourceLeaf{150}

\[
\begin{gathered}x=m'x'+\frac{m'N-m'b-n'c}{M}y',\\y=n'x'+\frac{n'N+m'a+n'b}{M}y'\end{gathered}
\]

in $(G)$ transit.\quad Vice versa, ex quavis transformatione propria formae $(F)$ in $(G)$ sequetur repraesentatio numeri $M$ per formam $(F)$, ad valorem $N$ pertinens. Scilicet si $(F)$ transit in $(G)$ positis $x=mx'-\nu y'$, $y=nx'+\mu y'$, $M$ repraesentatur per $(F)$ ponendo $x=m$, $y=n$, et quoniam hic $m\mu+n\nu=1$, valor expr. $\surd D(\operatorname{mod.}M)$ ad quem repraesentatio pertinet erit $\mu(bm+cn)-\nu(am+bn)$ \emph{i. e.} $N$.\quad Ex pluribus vero transformationibus propriis diversis, sequentur totidem repraesentationes diversae ad $N$ pertinentes\HistoricalNote{*)}{Si ex duabus transformationibus propriis diversis eadem repraesentatio defluere supponitur, illae ita se habere debebunt:
\[
\begin{aligned}
1)\quad &x=mx'-\nu y',\quad y=nx'+\mu y';\\
2)\quad &x=mx'-\nu'y',\quad y=nx'+\mu'x'.
\end{aligned}
\]
Sed ex duabus aequationibus
\[
\begin{gathered}
m\mu+n\nu=m\mu'+n\nu',\\
\mu(mb+nc)-\nu(ma+nb)=\mu'(mb+nc)-\nu'(ma+nb),
\end{gathered}
\]
facile deducitur esse aut $M=0$ aut $\mu=\mu'$, $\nu=\nu'$. At $M=0$ iam exclusimus.}. ---\quad Hinc facile colligitur, si omnes transformationes propriae formae $(F)$ in $(G)$ habeantur, ex his omnes repraesentationes ipsius $M$ per $(F)$ ad valorem $N$ pertinentes sequi.\quad Unde quaestio de repraesentationibus numeri dati per formam datam (in quibus indeterminatae valores inter se primos nanciscuntur) investigandis, reducta est ad quaestionem de inveniendis omnibus transformationibus propriis formae illius in datam aequivalentem.

Applicando iam ad haec ea quae in art. $162$ docuimus, facile concluditur: Si repraesentatio aliqua numeri $M$ per formam $(F)$ ad valorem $N$ pertinens sit haec: $x=\alpha$, $y=\gamma$: formulam generalem omnes repraesentationes eiusdem numeri per formam $(F)$, ad valorem $N$ pertinentes, comprehendentem fore hanc:

\[
x=\frac{\alpha t-(\alpha b+\gamma c)u}{m},\quad y=\frac{\gamma t+(\alpha a+\gamma b)u}{m}
\]

ubi $m$ divisor communis maximus numerorum $a,2b,c$; et $t,u$ omnes numeri, indefinite, aequationi $tt-Duu=mm$ satisfacientes.

\Article{170}

Si forma $(a,b,c)$ ancipiti alicui aequivalens, adeoque formae $(M,N,\frac{NN-D}{M})$ tam proprie, quam improprie, sive tam formae $(M,N,\frac{NN-D}{M})$, quam huic $(M,-N,\frac{NN-D}{M})$\quad proprie: repraesentationes numeri $M$ habebuntur per formam \SourceLeaf{151}$(F)$, tam ad valorem $N$, quam ad valorem $-N$, pertinentes.\quad Et vice versa si plures repraesentationes numeri $M$ per eandem formam $(F)$, ad valores \emph{oppositos} expr. $\surd D(\operatorname{mod.}M)$, $N,-N$, pertinentes habentur: forma $(F)$ formae $(G)$ tam proprie quam improprie aequivalens erit, formaque anceps assignari poterit, cui $(F)$ aequivaleat.

Haec generalia de repraesentationibus hic sufficiant: de repraesentationibus, in quibus indeterminatae valores inter se non primos habent, infra dicemus.\quad Respectu aliarum proprietatum, formae quarum determinans est negativus prorsus alio modo sunt tractandae, quam formae determinantis positivi: quare iam utrasque seorsim considerabimus.\quad Ab illis tamquam facilioribus initium facimus.

\begin{center}\emph{De formis determinantis negativi.}\end{center}

\Article{171}

\textsc{Problema.}\quad \emph{Proposita forma quacunque, $(a,b,a')$, cuius determinans negativus, $=-D$, designante $D$ numerum positivum, invenire formam huic proprie aequivalentem, $(A,B,C)$, in qua $A$ nec maior quam $\surd\tfrac43D,$ $C$, nec minor quam $2B$.}

\emph{Solutio.}\quad Supponimus in forma proposita non omnes tres conditiones simul locum habere: alioquin enim aliam formam quaerere opus non esset.\quad Sit $b'$ residuum abs. min. numeri $-b$ secundum modulum $a'$\HistoricalNote{*)}{Observare convenit, si formae alicuius $(a,b,a')$ terminus primus vel ultimus $a$ vel $a'$ sit $=0$, ipsius determinantem esse quadratum positivum: quare illud in casu praesenti evenire nequit. --- Ex simili ratione termini exteri $a,a'$ formae determinantis negativi, signa opposita habere non possunt.}, atque $a''=\frac{b'b'+D}{a'}$, qui erit integer quia $b'b'\equiv bb$, $b'b'+D\equiv bb+D\equiv aa'\equiv$ $0(\operatorname{mod.}a').$\quad Iam si $a''<a'$, fiat denuo $b''$ resid. abs. min. ipsius $-b'$ secundum mod. $a''$. atque $a'''=\frac{b''b''+D}{a''}$. Si hic iterum $a'''<a''$, sit rursus $b'''$ res. abs. min. ipsius $-b''$ secundum mod. $a'''$ atque $a''''=\frac{b'''b'''+D}{a'''}$.\quad Haec operatio continuetur donec in progressione $a',a''$, $a''',a''''$ etc. ad terminum $a^{m+1}$ perveniatur, qui praecedente suo $a^m$ non sit minor, quod tandem evenire debet, quia alias progressio infinita numerorum integrorum continuo decrescentium haberetur.\quad Tum forma $(a^m,b^m,a^{m+1})$ omnibus conditionibus satisfaciet.

\emph{Dem.}\quad I.\quad In progressione formarum $(a,b,a'),\ (a',b',a''),\ (a'',b'',a''')$ etc. quaevis praecedenti est contigua, quare ultima primae proprie aequivalens erit (artt. $159$, $160).$
